% Generated by scripts/assemble.py; corrections belong in source/corrections.json.
\examhead{182}{西北工业大学 2025 年数学分析试卷}{172}
% p173-b01; source PDF page 173
\blockanchor{p173-b01}
\question{1}(20 分) 用极限的严格数学定义证明:\par
% p173-b02; source PDF page 173
\blockanchor{p173-b02}
(1) $\lim_{x\to -\infty}\ln \left(1 + \frac{1}{x}\right) = 0.$\par
% p173-b03; source PDF page 173
\blockanchor{p173-b03}
(2) $\lim_{n\to \infty}\frac{e^n}{n!} = 0.$\par
% p173-b04; source PDF page 173
\blockanchor{p173-b04}
\question{2}(20 分) 解答如下问题:\par
% p173-b05; source PDF page 173
\blockanchor{p173-b05}
(1) 求极限 $\lim_{x\to 0^{+}}\frac{x^x - x^{\sin x}}{x^2}.$\par
% p173-b06; source PDF page 173
\blockanchor{p173-b06}
(2) 求级数的和 $\sum_{n=2}^{\infty} \frac{1}{n(n-1)3^n}$.\par
% p173-b07; source PDF page 173
\blockanchor{p173-b07}
\question{3}(15 分) 试用确界原理证明有限覆盖定理.\par
% p173-b08; source PDF page 173
\blockanchor{p173-b08}
\question{4}(15 分) 证明: 若 $f$ 是区间 $[a, b]$ 上只有有限个间断点的有界函数, 则 $f$ 在 $[a, b]$ 上黎曼可积.\par
% p173-b09; source PDF page 173
\blockanchor{p173-b09}
\question{5}(15 分) 计算 $\int_0^{+\infty} \frac{\cos ax - \cos bx}{x^2} \, \mathrm{d}x$, 其中 $b > a > 0$.\par
% p173-b10; source PDF page 173
\blockanchor{p173-b10}
\question{6}(15 分) 设 $n$ 为正整数, $x, y > 0$, 用条件极值方法证明 $\frac{x^n + y^n}{2} \geq \left(\frac{x + y}{2}\right)^n$.\par
% p173-b11; source PDF page 173
\blockanchor{p173-b11}
\question{7}(15 分) 设 $f(x)$ 在 $[0, +\infty)$ 上一致连续, 证明: 存在 $a, b > 0$ 使得 $|f(x)| \leq a|x| + b$.\par
% p173-b12; source PDF page 173
\par\Needspace{6\baselineskip}
\blockanchor{p173-b12}
\question{8}(15 分) 计算曲面积分\par
\nopagebreak[4]
% p173-b13; source PDF page 173
\blockanchor{p173-b13}
\displayformula{I = \iint_ {S} \frac {x \mathrm{d} y \mathrm{d} z + y \mathrm{d} z \mathrm{d} x + z \mathrm{d} x \mathrm{d} y}{\sqrt {(x ^ {2} + y ^ {2} + z ^ {2}) ^ {3}}}.}
% p173-b14; source PDF page 173
\blockanchor{p173-b14}
其中 S 为曲面  $\frac{x^{2}}{2} + \frac{y^{2}}{3} + \frac{z^{2}}{4} = 1$ ，取其外侧.\par
% p173-b15; source PDF page 173
\blockanchor{p173-b15}
\question{9}(15 分) 设 $f(x)$ 是定义在 $(- \infty, +\infty)$ 上的凸函数且有上界, 证明: $f(x)$ 是常数.\par
% p173-b16; source PDF page 173
\par\Needspace{6\baselineskip}
\blockanchor{p173-b16}
\question{10}(5 分) 设 $f(x)$ 在 $[a, b]$ 上连续, 非负且严格单调递增, 由积分中值定理, 对任意的正整数 $k$, 存在 $x_k \in [a, b]$, 使得\par
\nopagebreak[4]
% p173-b17; source PDF page 173
\blockanchor{p173-b17}
\displayformula{f ^ {k} (x _ {k}) = \frac {1}{b - a} \int_ {a} ^ {b} f ^ {k} (x) \mathrm{d} x.}
% p173-b18; source PDF page 173
\blockanchor{p173-b18}
证明： $\lim_{k\to \infty}x_k = b.$\par
